\section{Erdős \#287}

\textbf{Problem.} For any $k \geq 2$ and distinct integers $1 < n_1 < \cdots < n_k$ with
$\sum_{i=1}^{k} \frac{1}{n_i} = 1$, must we have $\max_i(n_{i+1} - n_i) \geq 3$?

\textbf{What was attempted.} We enumerated all representations of $1$ as a sum of distinct
unit fractions with denominators in $\{2, \ldots, N\}$ for $N \leq 30$, computed the
maximum consecutive gap $\max_i(n_{i+1}-n_i)$ for each representation, and searched for
any with maximum gap at most $2$ (a counterexample).

\textbf{What the computation showed.} After fixing two bugs---a \texttt{TypeError} from
incorrect list concatenation in round~1, and a pruning error that incorrectly returned
zero results in round~3---the corrected search (round~4) found 19 valid representations
for $N=20$ and confirmed all have maximum gap $\geq 3$. For example,
$(2,3,6)$ gives gaps $(1,3)$ with maximum gap $3$, saturating the bound.
No representation with all consecutive gaps $\leq 2$ was found within the search range.

\textbf{What went wrong or is inconclusive.} The search is limited to denominators up to
$N=30$ due to exponential blowup; the run at $N=40$ timed out after 60 seconds.
There is no reason the conjecture should fail only for small denominators, but the
search is far from exhaustive: representations using large denominators, or those
requiring many terms, are entirely unchecked. Finding no counterexample for small $N$
confirms the code works and is consistent with the conjecture, but constitutes no
mathematical progress. The conjecture remains open.

\textbf{Verdict:} No counterexample found for denominators up to $N=30$; result is
consistent with the conjecture but entirely inconclusive as a contribution toward
resolving it.